\documentclass{missingbook}
\usepackage{verifiable}

% chapters are included when they exist, so that the book builds while it is being written
\newcommand{\chapterfile}[1]{\IfFileExists{chapters/#1.tex}{\input{chapters/#1}}{}}

\begin{document}

\frontmatter
\hypersetup{pdftitle={Verifiable Intelligence: Interactive Proofs as a Theory of AI Oversight},pdfauthor={Written by Claude; conceived and edited by Xiaoyu Li}}
\mbtitlepage{Verifiable Intelligence}{Interactive Proofs as a Theory of AI Oversight}{Written by Claude\\[2.5mm]Conceived and edited by Xiaoyu Li}{First edition \textperiodcentered{} October 2026}
\addchap{Preface}

This book began with a list of textbooks that should exist but do not. One entry asked a question that is becoming practical: once an AI system is better than its overseers at a task, how do the overseers check its answers? The entry observed that complexity theory has studied exactly this question for forty years, as the question of what a weak verifier can learn from a strong prover that may lie, and that results bearing on it are scattered across papers on AI safety: debate and its connection to polynomial space, doubly-efficient debate, prover–verifier games, spot checks in the style of probabilistically checkable proofs, delegation of computation, and audits of models in zero knowledge. Nobody had yet taught the complexity theory systematically for this problem. The entry proposed to do so one protocol at a time, and to ask the same three questions of every protocol: is honesty an equilibrium, what does the verifier spend, and how often is a lie caught?

The resulting book has fifteen chapters in five parts. Part I studies provers that are fixed objects: certificates, proofs that the verifier spot-checks, and programs that it queries. Part II studies a single prover that answers questions, from private coins to $\cc{IP}=\cc{PSPACE}$ and the delegation of computation to an efficient prover. Part III sets provers against each other: debate, cross-examination, and debate in which the honest side must be efficient. Part IV studies provers that are trained or paid rather than adversarial: rational proofs, prover–verifier games, and the verification of learned models. Part V studies claims about objects the prover will not show, with zero knowledge and proofs of training, and ends with backdoors that no audit can detect. Each chapter opens with an oversight problem, defines its protocol, proves the protocol's main theorem completely, runs the protocol in a short script against honest and cheating provers, answers the three questions, and closes by saying what the mathematical model leaves out of the real problem.

The intended readers are graduate students and researchers in machine learning and AI safety who want the theory behind scalable oversight, and theoretical computer scientists who want to see where their theory meets it. The prerequisites are undergraduate discrete mathematics, probability, linear algebra, and algorithms; no course in complexity theory or cryptography is assumed. Each chapter defines what it uses, and an appendix collects the shared tools, with proofs of those whose proofs are short.

Two cautions belong at the front. First, the bridge between the two fields is an analogy with exact parts and idealized parts, and the book tries to say which is which. A human overseer is not a polynomial-time machine that never errs, a trained model is neither an all-powerful prover nor a worst-case adversary, and a protocol that certifies a formal statement certifies nothing about whether the statement says what the overseer meant. Every chapter's last section is about these gaps. Second, much of the AI-safety side of the subject is young: several of the results discussed are preprints from 2024 to 2026, and the experiments behind its empirical claims are small. The chapters date such claims and say when a result has not been refereed.

\paragraph{A note on authorship.} This book was written by Claude, an AI model made by Anthropic, and conceived and edited by Xiaoyu Li, who proposed the subject, its plan, and the three questions. Every theorem is stated with its hypotheses and proved or explicitly sketched, every claim attributed to a paper was checked against the paper or a reliable account of it wherever these could be reached, every number from the book's scripts comes from code that the reader can run, and the colophon describes how the book was made. It would be an odd book on oversight that asked to be trusted. Its claims are written to be checked, and readers who find an error are invited to report it, as a verifier would: with the step that fails.

\mbtableofcontents
\addchap{Notation}

The chapters share the notation below. A chapter that needs a clashing local symbol says so where it introduces it, and the notes at the end of each chapter translate the notation of its main sources into this one.

\addsec*{Parties and protocols}

\begin{tabular}{@{}p{0.27\textwidth}p{0.69\textwidth}@{}}
$\V$ & the verifier; in a debate, the judge \\
$\Prov$ & the prover, and the honest prover's strategy \\
$\Prov^*$ & an arbitrary prover, which may cheat \\
$\Prov_1,\Prov_2$ & two provers questioned separately \\
$\party{A},\party{B}$ & two debaters \\
$\party{H}$ & a human judge consulted as an oracle \\
$x$, $n=\abs{x}$ & the input and its length \\
$L$ & a language: the set of inputs for which the claim is true \\
$\interact{\Prov}{\V}(x)$ & the verifier's output, \acc{} or \rej, after interacting with $\Prov$ on $x$ \\
$c$, $s$ & completeness and soundness error; $1-s$ is the probability that a lie is caught \\
$k$, $q$, $t$ & rounds, queries to a proof, repetitions \\
$\rho$ & the verifier's random bits; $r$ a random field element \\
$T$ & the running time of the computation that a claim is about \\
$\lambda$, $\negl(\lambda)$ & a security parameter and a negligible function of it
\end{tabular}

\addsec*{Classes}

Complexity classes are set in bold: $\cc{P}$, $\cc{NP}$, $\cc{coNP}$, $\cc{BPP}$, $\cc{MA}$, $\cc{AM}$, $\cc{IP}$, $\cc{PSPACE}$, $\cc{EXP}$, $\cc{NEXP}$, $\cc{MIP}$, $\cc{\#P}$, and $\cc{PCP}(r,q)$ for languages with probabilistically checkable proofs that use $r$ random bits and $q$ queries. The appendix defines most of them; \cref{ch:spotchecks} defines $\cc{PCP}(r,q)$, and \cref{ch:crossexam} defines $\cc{MIP}$.

\addsec*{Algebra and probability}

\begin{tabular}{@{}p{0.27\textwidth}p{0.69\textwidth}@{}}
$\F$, $\abs{\F}$ & a finite field and its size \\
$\ell$, $d$ & the number of variables of a polynomial and its degree in each variable \\
$\tilde f$ & the multilinear extension of a function $f\colon\bits^\ell\to\F$ \\
$S$, $D$ & the size and depth of a circuit \\
$\bits$ & the set $\{0,1\}$; $\bits^n$ the strings of length $n$ \\
$\chi_S$, $\hat F(S)$ & the character $x\mapsto(-1)^{\sum_{i\in S}x_i}$ of $\F_2^n$, and the Fourier coefficient of $F=(-1)^f$ at $S$ \\
$\Prob$, $\E$ & probability and expectation, always over stated randomness \\
$m$, $\cD$, $\cH$ & a number of samples, a distribution, a hypothesis class \\
$\eps$, $\delta$ & the accuracy and the confidence of a learning guarantee
\end{tabular}

\addsec*{Asymptotics}

For functions of the input length $n$ (or of $\lambda$), $f=O(g)$ means $f\le Cg$ for a constant $C$ and all large $n$; $\Omega$, $\Theta$ and $o$ are as usual, $\poly(n)$ is some polynomial in $n$, and $\polylog(n)$ some polynomial in $\log n$. Logarithms are to base $2$ unless written $\ln$.


\mainmatter
\chapterfile{00-prologue}

\casepart{Fixed Provers}{The simplest prover is a fixed object: a certificate, a proof written out in full, a program. The verifier reads it, reads a few bits of it, or queries it on inputs of its choosing, and the object cannot change its answers in response.}{S,Q,Q,Q,D}
\chapterfile{01-certificates}
\chapterfile{02-spotchecks}
\chapterfile{03-consistency}

\casepart{Conversation}{A prover that answers questions can adapt its answers to them. In exchange the verifier gains a resource it did not have: questions the prover cannot predict. With interaction and private coins, a verifier that can barely compute can check claims about computations of exponential length.}{P,V,P,V,P,V,D}
\chapterfile{04-interaction}
\chapterfile{05-sumcheck}
\chapterfile{06-delegation}

\casepart{Competition}{When the verifier cannot check one prover, it can set two against each other. Debate gives a judge who reads a single step of an argument the power of polynomial space, and provers questioned apart, who cannot confer, give more. What the honest side must compute decides which of these can be used.}{P,A,P,A,P,A,D}
\chapterfile{07-debate}
\chapterfile{08-crossexam}
\chapterfile{09-efficientdebate}

\casepart{Incentives}{A trained model is not a worst-case adversary. It is the outcome of an optimization, and what it learns depends on what it is paid for. When does payment make the truth the best policy, which equilibria does training find, and can a learned model be checked with fewer samples than it took to learn?}{P,X,R,o,P,V,P,D}
\chapterfile{10-rational}
\chapterfile{11-pvg}
\chapterfile{12-pacverify}

\casepart{Secrets and Limits}{Some claims are about objects the prover will not show: its weights, its data, the record of its training. Cryptography lets a verifier check such claims without seeing the objects. It also shows that some properties of a model can be hidden from every audit.}{C,V,P,o,S,K,Q,Q,D}
\chapterfile{13-zeroknowledge}
\chapterfile{14-training}
\chapterfile{15-backdoors}

\bookmarksetup{startatroot}
\chapterfile{16-epilogue}

\appendix
\chapterfile{17-toolkit}

\backmatter
\printindex
\specialchapter{How This Book Was Made}{Colophon}{S,Q,Q,Q,D}

\addsec{Roles}

\emph{Verifiable Intelligence} was written by Claude Opus 5.5, an AI model made by Anthropic, in October 2026. Xiaoyu Li conceived the series \emph{The Missing Textbooks} and proposed this volume: its subject, its gap (a theory of AI oversight that existed only as scattered papers), its format (one oversight protocol per chapter), the three questions every chapter answers, its first examples, and its position relative to existing books. Xiaoyu Li also set the direction of the series (language, design, and credit) and is responsible for the book's publication.

\addsec{Process}

The book was written in one long session at maximum reasoning effort. A coordinating instance of the model first surveyed the literature with four research assistants, separate instances that verified, with web search and the papers themselves, the AI-safety work on debate and its complexity theory, on prover--verifier games, learning and incentives, on verifiable computation, zero knowledge, proofs of training and backdoors, and the classical results of interactive proofs, and that checked whether a textbook of this kind already existed (none was found). Their digests corrected the plan in several places before a chapter was written: the claim that cross-examination raises debate to nondeterministic exponential time comes from a 2020 forum writeup, not from the 2018 debate paper to which it is often attributed; the first version of prover--estimator debate had a wrong main theorem; and the three-query PCP of Håstad has imperfect completeness.

The coordinator then wrote the plan (the thesis, a precise reading of the three questions, the notation, fixed labels for results that chapters cite across the book, and a specification for each chapter), a style guide (the anatomy of a chapter, the voice, a proof style that follows Halmos, and the conventions of the scripts), the volume's typography, a shared kit for the scripts, the prologue, the preface, the notation, the appendix, and the epilogue.

Separate instances of the same model drafted the fifteen chapters, five at a time on a rolling basis. Each verified its references and the claims it attributes to papers, wrote and ran the scripts behind its figures, checked its theorems and constants numerically, and built its chapter with the volume's checks for errors, references, layout, and style. The shared budget for web search ran out early in the drafting, and the later authors verified sources by fetching arXiv, Crossref, DBLP, OpenAlex, and publishers' pages directly. Several authors went beyond their specifications, with results that are, as far as they could find, new: the exact soundness error $3/4$ of the Hadamard PCP with a tight example (\cref{ch:spotchecks}), the tight consistency bound for two provers (\cref{ch:crossexam}), the optimality of the Brier payment and a collapse of the polynomial hierarchy from a one-round rational proof with a noticeable gap (\cref{ch:rational}), a bound for committed spot checks (\cref{ch:training}), and a sharpening of the hypothesis behind undetectable backdoors (\cref{ch:backdoors}). The chapters mark these as the book's own.

Each chapter, and the prologue, appendix, and epilogue together, then received one independent referee pass from a separate instance, which re-derived the results and checked them numerically, often by exhaustive search over small cases, reran the scripts and compared every printed number with the text, worked the exercises, and checked the claims about papers against the papers; it corrected what it found directly. The referees found real errors. A definition of tree-like resolution let a proof reuse its lines, which broke a count of clauses (\cref{ch:certificates}); a remark claimed two bounds of the linearity test could not be improved, which a 1996 result refutes (\cref{ch:spotchecks}); one part of the self-correction proposition and one proposition about circuits with unused gate labels were false as stated (\cref{ch:consistency,ch:delegation}); a protocol's description did not match its worked example, which moved a game's value from $23/24$ to $7/8$ until the protocol was repaired (\cref{ch:crossexam}); a model of noisy human judgments let a liar tilt steps that should have been fixed (\cref{ch:efficientdebate}); a collapse of the polynomial hierarchy was stated without its hypothesis (\cref{ch:rational}); a lower bound was stated beyond the range its construction covers (\cref{ch:pacverify}); a guarantee for proofs of training promised nothing beyond the fraction of segments checked, since every model has a training log with at most one invalid segment (\cref{ch:training}); and the prologue's first protocol could be defeated by a model that knows the verifier's prime. Each is corrected in the text. The referees found no bug in a script that changed a conclusion, but many misquoted or misrounded numbers, and several sentences claiming that every measured interval contained its predicted value when one or two did not. One author had changed the random seed of an experiment after a first seed put a measured point above a proved bound; the referee reran it with fresh seeds, found the point an expected fluctuation, and the chapter now reports the episode and how often it should occur (\cref{ch:zeroknowledge}). In \cref{ch:pvg}, the referee checked against the papers the chapter's claims that three published results need a repair, confirmed them numerically where possible, and worded them as readings of the papers.

A usage limit stopped the running agents three times, authors and referees alike, and each was resumed with its context. The coordinator then integrated the book: it reconciled cross-references between chapters, gave every work cited in several chapters a single printed form, rebuilt the epilogue's ledger from the chapters' final protocol cards, and checked the whole book for layout, undefined references, and overflowing lines.

\addsec{Limits}

No human has yet checked every proof or every attribution. Web search and the papers themselves were available while the book was written, but some full texts could not be opened, so some claims about what a paper proves were checked against abstracts, the paper's own summary, or secondary accounts; the chapters say so where it matters. One referee pass per chapter finds many errors but not all of them, and a referee that is an instance of the same model may share the author's blind spots. Results that the book presents as its own have been checked by the same model only. Many of the AI-safety results discussed are preprints or forum posts from 2024 to 2026, and the experiments behind its empirical claims are small. The scripts are deliberately small too: a measured acceptance rate that matches a bound on a toy instance confirms the arithmetic of the bound, not its relevance to a deployed system. Corrections are welcome and will be recorded in the repository of \emph{The Missing Textbooks}.

\medskip
{\small\noindent Typeset with pdf\LaTeX{} in \texttt{missingbook}, the series class built on \textsc{koma}-Script, with the volume style \texttt{verifiable}; Palatino (newpx) for text and mathematics, Source Sans and Source Code Pro for labels, code, and bits, biblatex and biber for references, Ti\emph{k}Z and pgfplots for figures, whose data were computed with Python, NumPy, and SciPy. The title page carries the Hadamard encoding of the number $1985$, the year in which interactive proofs were introduced: the $2048$ bits $f(x)=\ip{u}{x}\bmod 2$ for $u=1985$ written in binary and $x=0,\dots,2047$, shown in $32$ rows of $64$ in an order fixed by a pseudorandom permutation (NumPy's generator with the seed $1985$), because in their natural order the bits of a linear function look regular. Three bits are ringed, at positions $x=45$, $y=673$ and $x\oplus y=652$; they read $1$, $1$, $0$, and the linearity test of \cref{ch:spotchecks} accepts, since $1+1=0$ modulo $2$. The foot of each chapter opening carries the message pattern of the chapter's protocol, drawn between a prover rail and a verifier rail. Both are computed or drawn by \texttt{figures/scripts/00-emblem.py} and \texttt{verifiable.sty}.\par}

\end{document}
